% Tutorial set: 02
% Source: Tutorial 2, Regression and classification, Questions 1--3
% Reference: Tutorial 2 reference solutions, slides 1--45
% Session date: 2026-08-26
% Human verified: Yes

\subsection{Tutorial 02：Regression、Logistic Classification 与 Feature Engineering}
\label{tutorial:SC4001-TUT02}

\exercisemeta
  {SC4001-TUT02}
  {2026-08-26}
  {Tutorial 2：Regression and classification，Questions 1--3；官方参考解答用于核对}
  {三题的模型、更新式与结果已核对；Q3 参考解答的系数不一致已订正}
  {已确认（2026-08-26）}

\exercisequestion{SC4001-TUT02-Q01}{Linear Neuron：SGD 与 Full-batch GD}

\textbf{对应讲义：}
\relatedtopic{SC4001-L02-T02}{Square-error Gradient、SGD 与 GD}；
\relatedtopic{SC4001-L02-T03}{Linear 与 Sigmoid Regression}；
\relatedtopic{SC4001-L02-T06}{Learning Rate 与 Autograd}

\subsubsection*{题目}

用下列八个 samples（样本）训练 three-input linear neuron（三输入线性神经元）$\hat y=\bm w^{\mathsf T}\bm x+b$：
\begin{center}
\small
\begin{tabular}{@{}rrrr@{}}
  \toprule
  $x_1$ & $x_2$ & $x_3$ & target $d$ \\
  \midrule
   0.09 & $-0.44$ & $-0.15$ & $-2.57$ \\
   0.69 & $-0.99$ & $-0.76$ & $-2.97$ \\
   0.34 &  0.65 & $-0.73$ &  0.96 \\
   0.15 &  0.78 & $-0.58$ &  1.04 \\
  $-0.63$ & $-0.78$ & $-0.56$ & $-3.21$ \\
   0.96 &  0.62 & $-0.66$ &  1.05 \\
   0.63 & $-0.45$ & $-0.14$ & $-2.39$ \\
   0.88 &  0.64 & $-0.33$ &  0.66 \\
  \bottomrule
\end{tabular}
\end{center}
初始 $\bm w=(0.77,0.02,0.63)^{\mathsf T}$、$b=0$、$\alpha=0.01$。分别展示一次 SGD（随机梯度下降）与 full-batch GD（全批量梯度下降）更新，并比较收敛曲线、最终 parameters（参数）及 predictions（预测值）。题目表头重复写了 $x_2$；第三维应为 $x_3$。

\begin{checkedsolution}
对 single sample（单个样本），令 $e=d-\hat y$，课程使用
\[
  \bm w\leftarrow\bm w+\alpha e\bm x,
  \qquad b\leftarrow b+\alpha e.
\]
参考解答 shuffle（打乱）后首先取 $\bm x=(0.34,0.65,-0.73)^{\mathsf T}$、$d=0.96$：
\[
  \hat y=-0.1851,\qquad e=1.1451,\qquad
  \bm w'=
  \begin{bmatrix}0.7739\\0.0274\\0.6216\end{bmatrix},\qquad
  b'=0.01145.
\]
SGD 每处理一个 sample 就更新一次；一个 epoch（训练轮次）共有八次更新，并应重新 shuffle samples。

对 full batch（完整批次），先用同一组旧参数计算
\[
  \hat{\bm y}=\bm X\bm w+b\bm 1,\qquad
  \bm e=\bm d-\hat{\bm y},
\]
再统一更新
\[
  \bm w\leftarrow\bm w+\alpha\bm X^{\mathsf T}\bm e,
  \qquad b\leftarrow b+\alpha\bm 1^{\mathsf T}\bm e.
\]
第一轮得到
\[
  \bm w'\approx
  \begin{bmatrix}0.7586\\0.1147\\0.6527\end{bmatrix},\qquad
  b'\approx-0.0736.
\]
该式使用 summed gradient（求和梯度）；若 loss implementation（损失实现）采用 mean reduction（均值归约），还会多出 $1/P$，需要相应调整 learning rate。

参考运行约在 200 epochs 后进入 plateau（平台期）：
\begin{center}
\small
\begin{tabular}{@{}lccc@{}}
  \toprule
  Method & $\bm w^{\mathsf T}$ & $b$ & MSE \\
  \midrule
  SGD & $(0.369,\ 2.566,\ -0.212)$ & $-1.165$ & 0.055 \\
  GD  & $(0.368,\ 2.567,\ -0.207)$ & $-1.163$ & 0.054 \\
  \bottomrule
\end{tabular}
\end{center}
两种方法的 training predictions（训练集预测）为：
\begin{center}
\small
\begin{tabular}{@{}crrr@{}}
  \toprule
  Sample & $d$ & $\hat y_{\mathrm{SGD}}$ & $\hat y_{\mathrm{GD}}$ \\
  \midrule
  1 & $-2.57$ & $-2.23$ & $-2.23$ \\
  2 & $-2.97$ & $-3.29$ & $-3.29$ \\
  3 &  0.96 &  0.78 &  0.78 \\
  4 &  1.04 &  1.01 &  1.01 \\
  5 & $-3.21$ & $-3.28$ & $-3.28$ \\
  6 &  1.05 &  0.92 &  0.92 \\
  7 & $-2.39$ & $-2.06$ & $-2.06$ \\
  8 &  0.66 &  0.87 &  0.88 \\
  \bottomrule
\end{tabular}
\end{center}
SGD 的具体轨迹依赖 shuffle order（打乱顺序）；上表是参考解答的一次运行结果，不是所有实现必须逐位复现的常数。
\end{checkedsolution}

\exercisequestion{SC4001-TUT02-Q02}{Logistic Neuron、BCE 与 3D Decision Boundary}

\textbf{对应讲义：}
\relatedtopic{SC4001-L02-T04}{Decision Boundary 与 Logistic Regression}；
\relatedtopic{SC4001-L02-T05}{Binary Cross-Entropy 与 Logistic Gradient}；
\relatedtopic{SC4001-L02-T06}{Learning Rate 与 Autograd}

\subsubsection*{题目}

使用 \texttt{make\_blobs} 生成 100 个 three-dimensional samples（三维样本），设置 two classes（两个类别）、$\texttt{cluster\_std}=5.0$。使用 \texttt{torch.nn.BCELoss()} 与 autograd（自动微分）训练 logistic neuron（逻辑神经元），报告收敛时的 classification error（分类错误）并画出 decision boundary（决策边界）。

\begin{checkedsolution}
参考数据设置为
\[
  \texttt{n\_samples=100},\quad
  \texttt{n\_features=3},\quad
  \texttt{centers=2},\quad
  \texttt{cluster\_std=5},\quad
  \texttt{random\_state=1}.
\]
模型与 BCE（binary cross-entropy，二元交叉熵）为
\[
  u=\bm w^{\mathsf T}\bm x+b,\qquad
  p=\sigma(u)=\frac{1}{1+e^{-u}},
\]
\[
  J=-\frac1P\sum_{p=1}^{P}
  \left[d_p\ln p_p+(1-d_p)\ln(1-p_p)\right].
\]
每个 epoch 依次执行 forward pass（前向传播）、\texttt{loss.backward()}、在 \texttt{torch.no\_grad()} 中更新 parameters，再将 gradients 清为 \texttt{None}。\texttt{BCELoss} 的输入是 sigmoid 后的 probabilities（概率）；参考代码把它命名为 \texttt{logits}，变量名并不准确。

使用 threshold $p>0.5$ 时，因 $\sigma(0)=0.5$ 且 sigmoid strictly increasing（严格递增），
\[
  p=0.5\iff u=0\iff\boxed{\bm w^{\mathsf T}\bm x+b=0}.
\]
故 decision boundary 在三维 feature space 中是一张 plane（平面）。参考运行报告
\[
  \bm w\approx(-0.0786,-0.3864,0.0437)^{\mathsf T},\qquad
  b\approx-0.0053,
\]
\[
  \mathrm{BCE}\approx0.306,\qquad
  \text{error count}=14,\qquad
  \text{error rate}=14\%.
\]
对应平面约为
\[
  -0.079x_1-0.386x_2+0.044x_3-0.005=0.
\]
数据的 \texttt{random\_state} 固定，但参考代码未固定 parameter-initialization
seed（参数初始化种子），所以不同运行的最终 weights 与 error 可能略有变化。

实践中常用更数值稳定的 \texttt{BCEWith\-Logits\-Loss}，但本题应按要求理解
\texttt{BCELoss} 与显式 sigmoid 的组合。
\end{checkedsolution}

\exercisequestion{SC4001-TUT02-Q03}{Scaled Sigmoid 与 Engineered Interaction Feature}

\textbf{对应讲义：}
\relatedtopic{SC4001-L02-T02}{Square-error Gradient、SGD 与 GD}；
\relatedtopic{SC4001-L02-T03}{Linear 与 Sigmoid Regression}；
\relatedtopic{SC4001-L02-T06}{Model Capacity 与 Autograd}

\subsubsection*{题目}

在 $0\le x,y\le1$ 上随机采样 100 个 training points（训练点），学习
\[
  \phi(x,y)=1.5+3.3x-2.5y+1.2xy.
\]
分别使用课程所称的 perceptron（连续 scaled-sigmoid neuron）与 linear neuron 训练，计算 training error（训练误差）、画出近似函数并比较两种方法。

\begin{checkedsolution}
在给定区域内
\[
  \frac{\partial\phi}{\partial x}=3.3+1.2y>0,
  \qquad
  \frac{\partial\phi}{\partial y}=-2.5+1.2x<0.
\]
因此
\[
  \phi_{\min}=\phi(0,1)=-1,\qquad
  \phi_{\max}=\phi(1,0)=4.8.
\]
参考解答据此使用 scaled sigmoid（缩放 Sigmoid）
\[
  \hat z=5.8\sigma(u)-1,\qquad
  u=w_1x+w_2y+b.
\]
严格地说，有限 $u$ 对应的 range（值域）是 $(-1,4.8)$，只能渐近接近两个端点。参考运行报告
\[
  \bm w\approx(2.95,-1.51)^{\mathsf T},\qquad
  b\approx-0.48,\qquad \mathrm{MSE}\approx0.014,
\]
即
\[
  \hat z=\frac{5.8}{1+\exp(0.48-2.95x+1.51y)}-1.
\]
单个 scaled-sigmoid neuron 只能表示 $f(w_1x+w_2y+b)$ 这一受限函数族，因此只能近似含 interaction term（交互项）$xy$ 的 bilinear surface（双线性曲面）。

Linear neuron 若只接收 raw features（原始特征）$(x,y)$，同样无法产生 $xy$ 项；但使用 feature mapping（特征映射）
\[
  \psi(x,y)=(x,y,xy)
\]
后，令 $x_1=x$、$x_2=y$、$x_3=xy$，即可写成
\[
  \boxed{\phi=\bm w^{\mathsf T}\psi(x,y)+b},\qquad
  \boxed{\bm w=(3.3,-2.5,1.2)^{\mathsf T},\quad b=1.5}.
\]
它在 transformed feature space（变换后的特征空间）中是 hyperplane（超平面），映射回原始 $(x,y)$ 空间则是 bilinear surface。

\begin{center}
\small
\begin{tabularx}{0.94\textwidth}{@{}>{\raggedright\arraybackslash}p{3.2cm}XX@{}}
  \toprule
  Model & Representation（表达能力） & 本题结论 \\
  \midrule
  Scaled-sigmoid neuron & $5.8\sigma(w_1x+w_2y+b)-1$ & 只能近似；参考 MSE $0.014$ \\
  Raw linear neuron & $w_1x+w_2y+b$ & 缺少 $xy$，不能精确表示 \\
  Engineered linear neuron & $w_1x+w_2y+w_3xy+b$ & 可用理论参数精确表示 \\
  \bottomrule
\end{tabularx}
\end{center}

\textbf{参考解答勘误：}题目及前半段解答均为 $1.2xy$，但后半段训练与比较页改成了 $0.2xy$，并给出 $w_3\approx0.29$。这些 numerical results（数值结果）对应另一目标函数，不能作为原题的正确 linear-neuron result；本解答以题目中的 $1.2xy$ 为准。
\end{checkedsolution}

\subsubsection*{本次 Tutorial 收口}

Q1 的核心是区分 per-sample update（逐样本更新）与 batch update（批量更新）；Q2 要把 sigmoid probability、BCE 与 $u=0$ 的 linear decision boundary 连起来；Q3 则说明 model linearity（模型线性）是相对于 features 而言，合适的 interaction feature 可以把原空间中的 nonlinear relation（非线性关系）变成特征空间中的线性关系。
